% GENERATED FILE. Edit canonical lessons, then run npm run generate:books.
% canonical-source-hash: fnv1a64:419475da581f54f1
% canonical-lessons: FR-C304-noyade, FR-C304-noyer, FR-C304-sauveteur, FR-C304-bouee, FR-C304-maitre-nageur

\chapter{Au bord de l'eau}
\label{ch:fr-au-bord-de-l-eau}
\hlchaptermodality{304}

\begin{chapteropening}
\textbf{By the end of this chapter:} \emph{I can talk about danger in the water: drowning, a rescuer, a lifebuoy and a lifeguard.}

Complete the last lesson of chapter 304: I can talk about danger in the water: drowning, a rescuer, a lifebuoy and a lifeguard.
\end{chapteropening}

\section[la noyade]{la noyade — drowning, a drowning}

\label{lesson:FR-C304-noyade}



\begin{quote}
Before the new one: say the French for to collapse, then the French for the rubble.
\end{quote}

\subsection*{You'll want to know: la noyade}
\textbf{la noyade} — \textquotedblleft{}drowning, a drowning\textquotedblright{}.

\begin{quote}
Il y a beaucoup de noyades en été.
\end{quote}

\begin{quote}
\emph{There are a lot of drownings in summer.}
\end{quote}

\begin{cousinweb}
\emph{noyade} is built on \emph{noyer} \textquotedblleft{}to drown\textquotedblright{} (EtymDB), which is inherited from Latin \emph{necāre} \textquotedblleft{}to kill\textquotedblright{}.
\end{cousinweb}

\begin{grammarlens}[title={What you've built}]
One more word for telling, explaining and describing at B1.
\end{grammarlens}

\subsection*{Your turn}
\begin{itemize}
\raggedright
  \item \emph{Say it:} \emph{la noyade}
  \item \emph{Say it:} \emph{la noyade}, once more
  \item \emph{Say it:} a sentence of your own with \emph{la noyade}
  \item \emph{Recall:} say the French for to collapse, then the French for the rubble, then say \emph{la noyade} again
\end{itemize}

\begin{tcolorbox}[breakable,colback=teal!4,colframe=teal!35!black,title={Before you move on}]
What does la noyade mean? (\textquotedblleft{}Drowning, a drowning\textquotedblright{}.) Say it once more.
\end{tcolorbox}
\section[noyer]{noyer — to drown (as se noyer); to flood}

\label{lesson:FR-C304-noyer}



\begin{quote}
Before the new one: say the French for the rubble, then the French for drowning.
\end{quote}

\subsection*{You'll want to know: noyer}
\textbf{noyer} — \textquotedblleft{}to drown (as se noyer); to flood\textquotedblright{}.

\begin{quote}
Il s'est presque noyé dans le lac.
\end{quote}

\begin{quote}
\emph{He nearly drowned in the lake.}
\end{quote}

Reflexive for a person who drowns: \emph{il s'est noyé}. Without \emph{se}, to drown someone or to flood: \emph{noyer le moteur}. The noun \emph{le noyer} is a walnut tree.

\begin{cousinweb}
\emph{noyer} is inherited from Latin \emph{necāre} \textquotedblleft{}to kill\textquotedblright{} (EtymDB), which in Late Latin narrowed to killing by water (Lewis \& Short).
\end{cousinweb}

\begin{grammarlens}[title={What you've built}]
One more word for telling, explaining and describing at B1.
\end{grammarlens}

\subsection*{Your turn}
\begin{itemize}
\raggedright
  \item \emph{Say it:} \emph{noyer}
  \item \emph{Say it:} \emph{noyer}, once more
  \item \emph{Say it:} a sentence of your own with \emph{noyer}
  \item \emph{Recall:} say the French for the rubble, then the French for drowning, then say \emph{noyer} again
\end{itemize}

\begin{tcolorbox}[breakable,colback=teal!4,colframe=teal!35!black,title={Before you move on}]
What does noyer mean? (\textquotedblleft{}To drown (as se noyer); to flood\textquotedblright{}.) Say it once more.
\end{tcolorbox}
\section[le sauveteur]{le sauveteur — the rescuer}

\label{lesson:FR-C304-sauveteur}



\begin{quote}
Before the new one: say the French for drowning, then the French for to drown.
\end{quote}

\subsection*{You'll want to know: le sauveteur}
\textbf{le sauveteur} — \textquotedblleft{}the rescuer\textquotedblright{}.

\begin{quote}
Les sauveteurs ont travaillé toute la nuit.
\end{quote}

\begin{quote}
\emph{The rescuers worked all night.}
\end{quote}

Feminine: \emph{la sauveteuse}.

\begin{cousinweb}
\emph{sauveteur} is built on \emph{sauver} \textquotedblleft{}to save\textquotedblright{}, which is inherited from Late Latin \emph{salvāre}, built on \emph{salvus} \textquotedblleft{}safe\textquotedblright{} (Lewis \& Short).
\end{cousinweb}

\begin{grammarlens}[title={What you've built}]
One more word for telling, explaining and describing at B1.
\end{grammarlens}

\subsection*{Your turn}
\begin{itemize}
\raggedright
  \item \emph{Say it:} \emph{le sauveteur}
  \item \emph{Say it:} \emph{le sauveteur}, once more
  \item \emph{Say it:} a sentence of your own with \emph{le sauveteur}
  \item \emph{Recall:} say the French for drowning, then the French for to drown, then say \emph{le sauveteur} again
\end{itemize}

\begin{tcolorbox}[breakable,colback=teal!4,colframe=teal!35!black,title={Before you move on}]
What does le sauveteur mean? (\textquotedblleft{}The rescuer\textquotedblright{}.) Say it once more.
\end{tcolorbox}
\section[la bouée]{la bouée — the buoy; the rubber ring}

\label{lesson:FR-C304-bouee}



\begin{quote}
Before the new one: say the French for to drown, then the French for the rescuer.
\end{quote}

\subsection*{You'll want to know: la bouée}
\textbf{la bouée} — \textquotedblleft{}the buoy; the rubber ring\textquotedblright{}.

\begin{quote}
On a lancé une bouée à l'homme dans l'eau.
\end{quote}

\begin{quote}
\emph{Someone threw a lifebuoy to the man in the water.}
\end{quote}

\emph{La bouée de sauvetage} — the lifebuoy; for children, \emph{une bouée} is a swim ring.

\begin{cousinweb}
\emph{bouée} is usually traced to Middle Dutch \emph{boeie} \textquotedblleft{}buoy\textquotedblright{}; the origin is not fully settled (EtymDB lists no parent).
\end{cousinweb}

\begin{grammarlens}[title={What you've built}]
One more word for telling, explaining and describing at B1.
\end{grammarlens}

\subsection*{Your turn}
\begin{itemize}
\raggedright
  \item \emph{Say it:} \emph{la bouée}
  \item \emph{Say it:} \emph{la bouée}, once more
  \item \emph{Say it:} a sentence of your own with \emph{la bouée}
  \item \emph{Recall:} say the French for to drown, then the French for the rescuer, then say \emph{la bouée} again
\end{itemize}

\begin{tcolorbox}[breakable,colback=teal!4,colframe=teal!35!black,title={Before you move on}]
What does la bouée mean? (\textquotedblleft{}The buoy; the rubber ring\textquotedblright{}.) Say it once more.
\end{tcolorbox}
\section[le maître-nageur]{le maître-nageur — the lifeguard, the swimming instructor}

\label{lesson:FR-C304-maitre-nageur}



\begin{quote}
Before the new one: say the French for the rescuer, then the French for the buoy.
\end{quote}

\subsection*{You'll want to know: le maître-nageur}
\textbf{le maître-nageur} — \textquotedblleft{}the lifeguard, the swimming instructor\textquotedblright{}.

\begin{quote}
Le maître-nageur surveille la piscine.
\end{quote}

\begin{quote}
\emph{The lifeguard watches over the pool.}
\end{quote}

Feminine: \emph{la maître-nageuse}.

\begin{cousinweb}
\emph{maître} is inherited from Latin \emph{magister} \textquotedblleft{}master\textquotedblright{}, and \emph{nageur} is built on \emph{nager} \textquotedblleft{}to swim\textquotedblright{}, which comes from Latin \emph{nāvigāre} \textquotedblleft{}to sail\textquotedblright{} (Lewis \& Short).
\end{cousinweb}

\begin{grammarlens}[title={What you've built}]
One more word for telling, explaining and describing at B1.
\end{grammarlens}

\subsection*{Your turn}
\begin{itemize}
\raggedright
  \item \emph{Say it:} \emph{le maître-nageur}
  \item \emph{Say it:} \emph{le maître-nageur}, once more
  \item \emph{Say it:} a sentence of your own with \emph{le maître-nageur}
  \item \emph{Recall:} say the French for the rescuer, then the French for the buoy, then say \emph{le maître-nageur} again
\end{itemize}

\begin{tcolorbox}[breakable,colback=teal!4,colframe=teal!35!black,title={Before you move on}]
What does le maître-nageur mean? (\textquotedblleft{}The lifeguard, the swimming instructor\textquotedblright{}.) Say it once more.
\end{tcolorbox}
